\documentclass[a4,11pt]{aleph-notas}

% -- Paquetes adicionales
\usepackage{enumitem}
\usepackage{aleph-comandos}

% -- Datos
\institucion{Facultad de Ciencias Exactas, Naturales y Ambientales}
\carrera{Catálogo STEM}
\asignatura{Matemática III}
\tema{Resumen 07: Campos vectoriales y regla de la cadena}
\autor[A. Merino]{Andrés Merino}
\fecha{Periodo 2026-1}

\logouno[0.16\textwidth]{Logos/logoPUCE_04_ac}
\definecolor{colortext}{HTML}{1957d1}
\definecolor{colordef}{HTML}{1957d1}
\fuente{montserrat}


\begin{document}

\encabezado

En este resumen consideraremos $n,m>1$, $\Omega\subset \R^n$ un conjunto abierto e $I\subset\R$ un intervalo.

%%%%%%%%%%%%%%%%%%%%%%%%%%%%%%%%%%%%%%
\section{Cálculo diferencial de campos vectoriales}
%%%%%%%%%%%%%%%%%%%%%%%%%%%%%%%%%%%%%%

\begin{defi}[Límite]
Dados $\func{F}{\Omega}{\R^m}$ un campo vectorial y $a\in \Omega'$, el límite de $F$ cuando $x$ tiende a $a$ es
\[
    \lim_{x\to a}F(x)
    =\left(\lim_{x\to a}f_1(x),\lim_{x\to a}f_2(x),\ldots,\lim_{x\to a}f_m(x)\right)
\]
siempre que los límites de la derecha existan.
\end{defi}

\begin{defi}[Continuidad]
Dados $\func{F}{\Omega}{\R^m}$ un campo vectorial y $a\in \Omega$, se dice que $F$ es continua en $a$ si
\[
    \lim_{x\to a}F(x)=F(a).
\]
Se dice que $F$ es continua en $\Omega$ si $F$ es continua en cada punto de $\Omega$.
\end{defi}

\begin{defi}[Derivada respecto a un vector]
    Dados $\func{F}{\Omega}{\R^m}$ un campo vectorial, $a\in \Omega$ y $y\in \R^n$, la derivada de $F$ en $a$ respecto a $y$ se define por
    \[
        F'(a;y)=\lim_{h\to 0}\frac{F(a+hy)-F(a)}{h},
    \]
    siempre que el límite exista.
\end{defi}

\begin{defi}[Derivadas parciales y matriz jacobiana]
    Dados $\func{F}{\Omega}{\R^m}$ un campo vectorial y $a\in \Omega$, la $k$-ésima derivada parcial de $F$ en $a$ es
    \[
        D_k F(a)=\frac{\partial F}{\partial x_k}(a)=F'(a;e^k).
    \]
    Si existen todas las derivadas parciales de $F$ en $a$, se define la matriz jacobiana de $F$ en $a$ por
    \[
        J_F(a) =
        \begin{bmatrix}
         D_1f_1(a) & D_2f_1(a) & \cdots & D_nf_1(a) \\
         D_1f_2(a) & D_2f_2(a) & \cdots & D_nf_2(a) \\
         \vdots & \vdots &  \ddots & \vdots\\
         D_1f_m(a) & D_2f_m(a) & \cdots & D_nf_m(a) \\
        \end{bmatrix}.
    \]
\end{defi}

\begin{prop}
    Sean $\func{F}{\Omega}{\R^m}$ un campo vectorial y $a\in \Omega$. Si existe la jacobiana de $F$ en $a$, entonces
    \[
        [F'(a;y)]= J_F(a) [y],
    \]
    para todo $y\in\R^n$.
\end{prop}

\begin{advertencia}
    La existencia de todas las derivadas direccionales no implica la continuidad del campo vectorial. Es decir, que un campo vectorial tenga jacobiana no implica que sea continuo.
\end{advertencia}

\begin{defi}[Diferenciabilidad]
    Dados $\func{F}{\Omega}{\R^m}$ un campo vectorial y $a\in \Omega$, se dice que $F$ es diferenciable en $a$ si existe una aplicación lineal $\func{T}{\R^n}{\R^m}$ tal que
    \[
        \lim_{h\to 0}\frac{F(a+h)-F(a)-T(h)}{\|h\|}=0.
    \]
    A la aplicación $T$ se la llama el diferencial de $F$ en $a$ y se lo denota por $F'(a)$ o $DF(a)$.
\end{defi}

\begin{teo}
    Sean $\func{F}{\Omega}{\R^m}$ un campo vectorial y $a\in \Omega$. Si $F$ es diferenciable en $a$ entonces
    \[
        [F'(a)]=J_F(a),
    \]
    es decir, $[F'(a)(h)]=J_F(a) [h]$ para todo $h\in\R^n$.
\end{teo}

\begin{teo}
    Sean $\func{F}{\Omega}{\R^m}$ un campo vectorial y $a\in \Omega$. Si $F$ es diferenciable en $a$, entonces $F$ es continuo en $a$.
\end{teo}

\begin{advertencia}
    Si un campo vectorial tiene jacobiana, no necesariamente es diferenciable.
\end{advertencia}

%%%%%%%%%%%%%%%%%%%%%%%%%%%%%%%%%%%%%%
\section{Regla de la cadena}
%%%%%%%%%%%%%%%%%%%%%%%%%%%%%%%%%%%%%%

En esta sección tomaremos $k\in\N^*$ y $\Omega_1\subset \R^k$, $\Omega_2\subset \R^n$ conjuntos abiertos.

\begin{teo}[Regla de la cadena]
    Sean $\func{F}{\Omega_1}{\R^m}$ y $\func{G}{\Omega_2}{\R^k}$ campos vectoriales tales que $\img(G)\subset \Omega_1$ y $a\in \Omega_2$. Si $G$ es diferenciable en $a$ y $F$ es diferenciable en $G(a)$, entonces $F\circ G$ es diferenciable en $a$ y se tiene que
    \[
        (F\circ G)'(a) = F'(G(a))\circ G'(a).
    \]
\end{teo}

\begin{prop}
    Bajo las hipótesis del teorema anterior,
    \[
        J_{F\circ G}(a) = J_F(G(a))\, J_G(a),
    \]
    así, si $i=1,\ldots,n$ y $j=1,\ldots,m$, entonces
    \[
        D_i (F\circ G)_j(a) = D_1 f_j(b)\, D_i g_1(a) + D_2 f_j(b)\, D_i g_2(a) + \cdots + D_k f_j(b)\, D_i g_k(a),
    \]
    donde $b=G(a)$.
\end{prop}

\begin{advertencia}
    Si se tiene $b=x(a)$ con
    \begin{gather*}
        \funcion{x}{\Omega_2\subset \R^n}{\R^k}{(t_1,\ldots,t_n)}{x(t_1,\ldots,t_n),}\\
        \funcion{y}{\Omega_1\subset \R^k}{\R^m}{(x_1,\ldots,x_k)}{y(x_1,\ldots,x_k),}
    \end{gather*}
    entonces
    \[
        \frac{\partial (y\circ x)_j}{\partial t_i}(a)
            = \frac{\partial y_j}{\partial x_1}(b)\,\frac{\partial x_1}{\partial t_i}(a)
            + \frac{\partial y_j}{\partial x_2}(b)\,\frac{\partial x_2}{\partial t_i}(a)
            + \cdots
            + \frac{\partial y_j}{\partial x_k}(b)\,\frac{\partial x_k}{\partial t_i}(a).
    \]
\end{advertencia}


\end{document}
